\writeSectionFrame{\presentationTitle}{Dependency Injection}
\begin{frame}[fragile]{Mit DI - Beispiel Strategie-Muster}
	\begin{exampleblock}{Monster.java}
		\begin{lstlisting}
public class Monster {
	private AttackStrategy attackStrategy;
	
	public Monster(AttackStrategy attackStrategy) {
		this.attackStrategy = attackStrategy;
	}
	
	public void changeAttackStrategy(AttackStrategy attackStrategy) {
		this.attackStrategy = attackStrategy;
	}
	
	public void attack() {
		attackStrategy.attack();
	}
}
		\end{lstlisting}
	\end{exampleblock}
\end{frame}
\note{
	Im Konstruktor übergeben wir von außen ein konkretes Objekt. Unser
	Monster weiß nicht mehr, mit welcher Strategie es genau arbeitet. Und
	ihm ist es egal, wo diese Implementierung herkommt. Wir brauchen auch 
	für den Strategiewechsel keine unterschiedlichen Methoden mehr. Wenn jetzt
	neue Strategien entwickelt werden, müssen wir Monster nicht mehr verändern.
}


\begin{frame}[fragile]{Verwendung von DI}
	\begin{exampleblock}{Main.java}
		\begin{lstlisting}
Monster fireMonster = new Monster(new FireAttackStrategy());
Monster meleeMonster = new Monster(new MeleeAttackStrategy());
Monster monster = new Monster(new MeleeAttackStrategy());

fireMonster.attack();  
meleeMonster.attack(); 
monster.attack();
System.out.println("Monster findet ein Schwert");
monster.changeAttackStrategy(new CloseCombatStrategy());
monster.attack();
		\end{lstlisting}
	\end{exampleblock}
\end{frame}

\begin{frame}{Vorteile von Dependency Injection}
	\begin{itemize}
	    \item Monster kennt keine konkreten Strategien mehr $\rightarrow$ lose Kopplung
	    \item Neue Strategien können hinzugefügt werden, ohne Monster zu ändern
	    \item Einfache Testbarkeit durch Mock- oder Teststrategien
	    \item Unterstützt flexible Konfiguration (zur Laufzeit oder über Frameworks)
	\end{itemize}
\end{frame}

\begin{frame}{Wichtige Arten der Dependency Injection}
	\begin{itemize}
	    \item \textbf{Constructor Injection:} Abhängigkeit wird im Konstruktor übergeben (wie im Beispiel)
	    \item \textbf{Setter Injection:} Abhängigkeit wird über Setter-Methoden gesetzt
	\end{itemize}
\end{frame}

\begin{frame}{Konstruktor-Injection} 
    \begin{exampleblock}{Vorteile} 
        \begin{itemize} 
            \item Fördert Unveränderlichkeit 
            \item Klare Abhängigkeiten 
            \item Einfache Testbarkeit 
            \item Zwingt Bereitstellung von Abhängigkeiten
        \end{itemize} 
    \end{exampleblock}
    \begin{alertblock}{Nachteile}
        \begin{itemize} 
            \item Umfangreicher Konstruktor bei vielen Abhängigkeiten 
            \item Probleme bei zirkulären Abhängigkeiten
        \end{itemize} 
    \end{alertblock}
\end{frame}

\note{ 
    Konstruktor-Injection fördert unveränderbare Objekte, da Abhängigkeiten 
    direkt gesetzt werden und nach der Erstellung des Objekts nicht mehr 
    geändert werden können. Der Konstruktor macht Abhängigkeiten deutlich 
    sichtbar, was den Code lesbarer und wartbarer macht. Bei Tests ist es 
    einfacher, Mock-Objekte oder alternative Implementierungen bereitzustellen.

    Ein Nachteil dieser Methode ist, dass Konstruktoren unübersichtlich 
    werden können, wenn viele Abhängigkeiten existieren. Außerdem können 
    zirkuläre Abhängigkeiten nicht ohne weiteres aufgelöst werden, was zu 
    Problemen führen kann. 
}

\begin{frame}{Setter-Injection} 
    \begin{exampleblock}{Vorteile} 
        \begin{itemize} 
            \item Flexibilität bei Abhängigkeitsänderung 
            \item Gut für optionale Abhängigkeiten 
            \item Löst zirkuläre Abhängigkeiten 
        \end{itemize} 
    \end{exampleblock}
    \begin{alertblock}{Nachteile}
        \begin{itemize} 
            \item Abhängigkeiten weniger offensichtlich 
            \item Mutierbarkeit des Objekts 
            \item Risiko von fehlenden Abhängigkeiten 
        \end{itemize} 
    \end{alertblock}       
\end{frame}

\note{ 
    Setter-Injection bietet Flexibilität, da Abhängigkeiten nach der 
    Objekterstellung geändert werden können. Sie eignet sich besonders gut, 
    wenn einige Abhängigkeiten optional sind oder wenn man mit zirkulären 
    Abhängigkeiten umgehen muss.

    Der Nachteil ist, dass Abhängigkeiten weniger offensichtlich sind, da sie 
    nicht sofort im Konstruktor ersichtlich sind. Außerdem bleibt das Objekt 
    mutierbar, was potenziell zu Fehlern führen kann. Ein weiteres Risiko besteht 
    darin, dass der Entwickler vergessen könnte, eine notwendige Abhängigkeit zu 
    setzen. 
}

\begin{frame}{Testen}
	\begin{block}{Frage}
		Wie hilft DI beim Testen?
	\end{block}
\end{frame}

\begin{frame}[fragile]{Testen ohne DI}
	\begin{exampleblock}{ExceptionAttackStrategy.java}
		\begin{lstlisting}
public class ExceptionAttackStrategy implements AttackStrategy {
	@Override
	public void attack() {
		// Simuliert einen Fehler in der Strategie
		throw new IllegalStateException("Fehlerhaft");
	}
}
		\end{lstlisting}
	\end{exampleblock}
\end{frame}

\begin{frame}[fragile]{Testen ohne DI}
	\begin{exampleblock}{NotDiMonster.java}
		\begin{lstlisting}
public class NotDiMonster {
	private final AttackStrategy attackStrategy;
	private int posX;
	private int posY;
	
	public NotDiMonster() {
		this.attackStrategy = new ExceptionAttackStrategy();
		posX = 0;
		posY = 0;
	}
	
	public void moveAndAttack() {
		posX++;
		posY++;
		attackStrategy.attack();
	}
}
		\end{lstlisting}
	\end{exampleblock}
\end{frame}

\begin{frame}[fragile]{Testen ohne DI}
	\begin{exampleblock}{MonsterTest.java}
		\begin{lstlisting}
@Test
void testMoveWithDI() {
	NotDiMonster monster = new NotDiMonster();
	monster.moveAndAttack();
	assertEquals(1, monster.getPosX());
	assertEquals(1, monster.getPosY());
}
		\end{lstlisting}
	\end{exampleblock}
	\pause
	\begin{alertblock}{Test schlägt fehl}
		Test schlägt fehl, weil AttackStrategy fehlerhaft ist, aber der Fehler
		ist nicht in der Klasse, die wir eigentlich testen möchten.
	\end{alertblock}
\end{frame}

\begin{frame}[fragile]{Verwendung von DI}
	\begin{exampleblock}{Monster.java}
		\begin{lstlisting}
public class Monster {
	private final AttackStrategy attackStrategy;
	private int posX;
	private int posY;
	
	public Monster(AttackStrategy attackStrategy) {
		this.attackStrategy = attackStrategy;
		posX = 0;
		posY = 0;
	}
	
	public void moveAndAttack() {
		posX++;
		posY++;
		attackStrategy.attack();
	}
}
		\end{lstlisting}
	\end{exampleblock}
\end{frame}

\begin{frame}{Testen}
	\begin{alertblock}{Problem}
		Zum Testen benötige ich eine konkrete Strategie. Aber was ist,
		wenn die Strategie-Klasse einen Fehler enthält?
	\end{alertblock}
	\pause
	\begin{exampleblock}{Mocking}
		Wir mocken die Strategie, d.h. wir bauen eine Implementierung, die 
		auf jeden Fall funktioniert. So kann Monster unabhängig von der 
		Angriffsstrategie getestet werden.
	\end{exampleblock}
	\pause
	\begin{block}{Hinweis}
		Wir implementieren Mocking hier ganz einfach. In professionellen
		Projekte würde man ein Mocking-Framework wie z.B. Mockito verwenden.
		Hier geht es nur ums Verständnis.
	\end{block}
\end{frame}

\begin{frame}[fragile]{Verwendung von DI}
	\begin{exampleblock}{MockedAttackStrategy.java}
		\begin{lstlisting}
public class MockedAttackStrategy 
		implements AttackStrategy {

	@Override
	public void attack() {
		System.out.println("...");
	}
}
		\end{lstlisting}
	\end{exampleblock}
\end{frame}

\begin{frame}[fragile]{Verwendung von DI}
	\begin{exampleblock}{MonsterTest.java}
		\begin{lstlisting}
@Test
void testMove() {
	AttackStrategy strategy = new MockedAttackStrategy();
	Monster monster = new Monster(strategy);
	monster.moveAndAttack();
	assertEquals(1, monster.getPosX());
	assertEquals(1, monster.getPosY());
}
		\end{lstlisting}
	\end{exampleblock}
\end{frame}

\begin{frame}{DI verbessert Testbarkeit}
	\begin{columns}[T,onlytextwidth]
	\column{0.48\textwidth}
		\textbf{Ohne DI}
		\begin{itemize}
		    \item Monster erstellt selbst die Strategie
		    \item Test schlägt fehl, wenn Strategie fehlerhaft ist
		    \item Test isoliert Monster nicht → Kopplung
		\end{itemize}
		
		\begin{exampleblock}{Beispiel}
			NotDiMonster → ExceptionAttackStrategy
		\end{exampleblock}
		
		\column{0.48\textwidth}
		\textbf{Mit DI}
		\begin{itemize}
		    \item Monster erhält Strategie von außen
		    \item Test kann sichere Mock-Strategie verwenden
		    \item Monster unabhängig von konkreter Strategie testbar
		\end{itemize}
		
		\begin{exampleblock}{Beispiel}
			Monster → MockedAttackStrategy
		\end{exampleblock}
	\end{columns}
\end{frame}
\note{
	Durch Entkopplung von Abhängigkeiten können wir Klassen isoliert testen. 
	Wir geben Monster einfach eine sichere Strategie, die garantiert keinen 
	Fehler enthält, und können so die Logik der Monster-Klasse testen, ohne 
	von der Strategie abhängig zu sein.
}

\begin{frame}{Weitere Möglichkeiten}
	\begin{itemize}
	  \item DI trennt Monster von Strategie
	  \item Man kann also einbauen:
	  	\begin{itemize}
	  	  \item fehlerhafte Strategien
	  	  \item langsame Strategien
	  	  \item schnelle Strategien
	  	  \item komplexe Strategien
	  	  \item \ldots
	  	\end{itemize}
	\end{itemize}
\end{frame}

\begin{frame}{Erweiterungen}
	\begin{itemize}
	  \item Monster und Angriffsstrategien benötigen weitere Konfiguration
	  \item eine Fabrik für Strategien
	  \item eine Fabrik für Monster
	  \item eine Konfiguration der Fabriken als Singleton
	  \item natürlich alles über DI
	  \item Verdrahtung der Objekte an zentraler Stelle
	  \item möglichst lose Kopplung
	\end{itemize}
\end{frame}

\begin{frame}[fragile]{Fabrik für die Angriffsstrategien}
	\begin{exampleblock}{AttackStrategyFactory.java}
		\begin{lstlisting}
public interface AttackStrategyFactory {
	AttackStrategy createAttackStrategy();
}
		\end{lstlisting}
	\end{exampleblock}
\end{frame}

\begin{frame}[fragile]{Fabrik für die Angriffsstrategien}
	\begin{exampleblock}{AttackStrategyFactory.java}
		\begin{lstlisting}
public class FireStrategyFactory implements AttackStrategyFactory {
	@Override
	public AttackStrategy createAttackStrategy() {
		AttackStrategy attackStrategy = new FireAttackStrategy();
		// Weitere Konfiguration der Strategie
		return attackStrategy;
	}
}
		\end{lstlisting}
	\end{exampleblock}
\end{frame}

\begin{frame}[fragile]{Fabrik für die Angriffsstrategien}
	\begin{exampleblock}{MonsterFactory.java}
		\begin{lstlisting}
public abstract class MonsterFactory {
	private AttackStrategyFactory attackStrategyFactory;
	
	public MonsterFactory(AttackStrategyFactory 
			attackStrategyFactory) {
		this.attackStrategyFactory = attackStrategyFactory;
	}
	
	public abstract Monster createMonster();
	
	protected AttackStrategyFactory getAttackStrategyFactory() {
		return attackStrategyFactory;
	}
}
		\end{lstlisting}
	\end{exampleblock}
\end{frame}

\begin{frame}[fragile]{Fabrik für die Angriffsstrategien}
	\begin{exampleblock}{GoblinFactory.java}
		\begin{lstlisting}
public class GoblinFactory extends MonsterFactory {
	public GoblinFactory(AttackStrategyFactory 
			attackStrategyFactory) {
		super(attackStrategyFactory);
	}

	@Override
	public Monster createMonster() {
		AttackStrategy strategy = getAttackStrategyFactory().
			createAttackStrategy();
		Monster monster = new Goblin(strategy);
		// Boesen Goblin weiter konfigurieren.
		return monster;
	}
}
		\end{lstlisting}
	\end{exampleblock}
\end{frame}

\begin{frame}[fragile]{Fabrik für die Angriffsstrategien}
	\begin{exampleblock}{FactoryConfiguration.java}
		\begin{lstlisting}
public class FactoryConfiguration {
	private static FactoryConfiguration instance;
	private final MonsterFactory evilDwarfFactory;
	private final AttackStrategyFactory meleeStrategyFactory;
	
	private FactoryConfiguration() {
		this.meleeStrategyFactory = new MeleeStrategyFactory();
		this.evilDwarfFactory = new 
			EvilDwarfFactory(meleeStrategyFactory);
	}
	
	public static FactoryConfiguration getInstance(){
		if (instance == null) {
			instance = new FactoryConfiguration();
		}
		
		return instance;
	}
}
		\end{lstlisting}
	\end{exampleblock}
\end{frame}

